\par\Needspace{6\baselineskip}\medskip
\SourceLeaf{196}\begin{center}\hypertarget{native-196-article145}{}\pdfbookmark[2]{145.}{article-145}145.\end{center}
Sind die Orte des Himmelskörpers auf seiner Bahn auf diese Weise bestimmt, so beginnt man mit zwei Berechnungen der Bahnelemente. Für die eine verbindet man den zweiten Ort mit dem dritten, für die andere den ersten mit dem zweiten und verwendet jeweils das zugehörige Zeitintervall. Bevor man damit beginnt, müssen jedoch die Zeitintervalle selbst korrigiert werden, wenn man die Aberration nach dem dritten Verfahren aus Artikel 118 berücksichtigen will. In diesem Fall sind die wirklichen Zeiten durch fiktive zu ersetzen, die jeweils um $493\rho,493\rho',493\rho''$ Sekunden früher liegen. Zur Berechnung der Abstände $\rho,\rho',\rho''$ dienen die Formeln
\[
\begin{aligned}
\rho&=\frac{R\sin(AD'-\zeta)}{\sin(\zeta-AD'+\delta)}=\frac{r\sin(AD'-\zeta)}{\sin\delta}\\[3pt]
\rho'&=\frac{R'\sin(\delta'-z)}{\sin z}=\frac{r'\sin(\delta'-z)}{\sin\delta'}\\[3pt]
\rho''&=\frac{R''\sin(A''D'-\zeta'')}{\sin(\zeta''-A''D'+\delta'')}=\frac{r''\sin(A''D'-\zeta'')}{\sin\delta''}\mathpunct{\cdot}
\end{aligned}
\]

Wären die Beobachtungen dagegen schon zu Beginn nach dem ersten oder zweiten Verfahren aus Artikel 118 um die Aberration berichtigt worden, so wäre diese Rechnung wegzulassen. Dann müssten auch die Abstände $\rho,\rho',\rho''$ nicht bestimmt werden, außer vielleicht zur Kontrolle, ob die der Berechnung der Aberrationen zugrunde gelegten Werte hinreichend genau waren. Schließlich entfällt die ganze Rechnung selbstverständlich auch dann, wenn man die Aberration gänzlich außer Betracht lassen will.
